\begin{exposition}
Zorich's discussion of approximate numerical computation in
\cite{ZorichMathematicalAnalysisI} is more than an applied aside. It is a
finite version of the same error bookkeeping that later becomes sequence
limits, Landau notation, and differentiability. An exact quantity is written
as an approximation plus an error:
\[
  x=\widetilde{x}+\alpha.
\]
In limits the same pattern is written
\[
  x_n=A+\alpha_n,\qquad \alpha_n\to 0.
\]
In differential calculus it becomes
\[
  f(a+h)=f(a)+Lh+r(h),\qquad r(h)=o(h).
\]
The names change, but the mechanism is the same: isolate the main term,
track the error, and prove that the remaining term is controlled at the
required scale.
\end{exposition}

\begin{definitionbox}{Definition (Absolute Error of an Approximation)}
\begin{definition}[Absolute Error of an Approximation]
\label{def:absolute-error-approximation}
Let $x\in\mathbb{R}$ be an exact value and let $\widetilde{x}\in\mathbb{R}$
be an approximation to $x$. The \textbf{absolute error} of the approximation
is
\[
  \Delta(\widetilde{x}) := |x-\widetilde{x}|.
\]
\end{definition}
\end{definitionbox}

\begin{remark*}[Standard quantified statement]
\[
\forall x,\widetilde{x}\in\mathbb{R},\quad
\Delta(\widetilde{x})=|x-\widetilde{x}|.
\]
\end{remark*}

\begin{remark*}[Predicate reading]
\[
\Delta(\widetilde{x})=|x-\widetilde{x}|.
\]
\end{remark*}

\begin{remark*}[Interpretation]
Absolute error measures the raw distance between an exact quantity and the value used in its place.
\end{remark*}

\begin{dependencies}
\begin{itemize}
  \item \hyperref[def:absolute-value]{Absolute Value}
\end{itemize}
\end{dependencies}

\begin{definitionbox}{Definition (Relative Error of an Approximation)}
\begin{definition}[Relative Error of an Approximation]
\label{def:relative-error-approximation}
Let $x\in\mathbb{R}$ be an exact value and let
$\widetilde{x}\in\mathbb{R}\setminus\{0\}$ be an approximation to $x$. The
\textbf{relative error} of the approximation is
\[
  \delta(\widetilde{x})
  :=
  \frac{\Delta(\widetilde{x})}{|\widetilde{x}|}.
\]
\end{definition}
\end{definitionbox}

\begin{remark*}[Standard quantified statement]
\[
\forall x\in\mathbb{R}\;\forall \widetilde{x}\in\mathbb{R}\setminus\{0\},\quad
\delta(\widetilde{x})=\frac{\Delta(\widetilde{x})}{|\widetilde{x}|}.
\]
\end{remark*}

\begin{remark*}[Predicate reading]
\[
\delta(\widetilde{x})=\Delta(\widetilde{x})/|\widetilde{x}|.
\]
\end{remark*}

\begin{remark*}[Interpretation]
In practice the exact value $x$ is usually unknown, so the exact value of
$\Delta(\widetilde{x})$ is unknown as well. What matters for analysis is an
upper bound: for example, $\Delta(\widetilde{x})\leq E$. The whole section is
about how such bounds behave under algebraic operations.
\end{remark*}

\begin{dependencies}
\begin{itemize}
  \item \hyperref[def:absolute-error-approximation]{Absolute Error of an Approximation}
  \item \hyperref[def:absolute-value]{Absolute Value}
\end{itemize}
\end{dependencies}

\begin{theorembox}{Theorem (Absolute Error of a Sum)}
\begin{theorem}[Absolute Error of a Sum]
\label{thm:absolute-error-sum}
Let $x,y,\widetilde{x},\widetilde{y}\in\mathbb{R}$. Then
\[
  \Delta(\widetilde{x}+\widetilde{y})
  \leq
  \Delta(\widetilde{x})+\Delta(\widetilde{y}).
\]

\smallskip\noindent
\hyperref[prf:absolute-error-sum]{\textit{Go to proof.}}
\end{theorem}
\end{theorembox}

\begin{remark*}[Standard quantified statement]
\[
\forall x,y,\widetilde{x},\widetilde{y}\in\mathbb{R},\quad
\Delta(\widetilde{x}+\widetilde{y})
\leq \Delta(\widetilde{x})+\Delta(\widetilde{y}).
\]
\end{remark*}

\begin{remark*}[Predicate reading]
\[
\Delta(\widetilde{x}+\widetilde{y})\leq \Delta(\widetilde{x})+\Delta(\widetilde{y}).
\]
\end{remark*}

\begin{remark*}[Interpretation]
Absolute error is subadditive under addition: the worst-case error in a sum is bounded by the sum of the component errors.
\end{remark*}

\begin{dependencies}
\begin{itemize}
  \item \hyperref[def:absolute-error-approximation]{Absolute Error of an Approximation}
  \item \hyperref[thm:absolute-value-sum-bound]{Generalised Triangle Inequality}
\end{itemize}
\end{dependencies}

\begin{theorembox}{Theorem (Absolute Error of a Product)}
\begin{theorem}[Absolute Error of a Product]
\label{thm:absolute-error-product}
Let $x,y,\widetilde{x},\widetilde{y}\in\mathbb{R}$. Then
\[
  \Delta(\widetilde{x}\widetilde{y})
  \leq
  |\widetilde{x}|\Delta(\widetilde{y})
  +
  |\widetilde{y}|\Delta(\widetilde{x})
  +
  \Delta(\widetilde{x})\Delta(\widetilde{y}).
\]

\smallskip\noindent
\hyperref[prf:absolute-error-product]{\textit{Go to proof.}}
\end{theorem}
\end{theorembox}

\begin{remark*}[Standard quantified statement]
\[
\forall x,y,\widetilde{x},\widetilde{y}\in\mathbb{R},\quad
\Delta(\widetilde{x}\widetilde{y})
\leq |\widetilde{x}|\Delta(\widetilde{y})
+|\widetilde{y}|\Delta(\widetilde{x})
+\Delta(\widetilde{x})\Delta(\widetilde{y}).
\]
\end{remark*}

\begin{remark*}[Predicate reading]
\[
\Delta(\widetilde{x}\widetilde{y})\leq |\widetilde{x}|\Delta(\widetilde{y})+|\widetilde{y}|\Delta(\widetilde{x})+\Delta(\widetilde{x})\Delta(\widetilde{y}).
\]
\end{remark*}

\begin{remark*}[Interpretation]
The two terms
$|\widetilde{x}|\Delta(\widetilde{y})$ and
$|\widetilde{y}|\Delta(\widetilde{x})$ are first-order: each contains one
error factor. The term
$\Delta(\widetilde{x})\Delta(\widetilde{y})$ is second-order: it contains
two error factors. This is the finite precursor of little-o bookkeeping.
\end{remark*}

\begin{dependencies}
\begin{itemize}
  \item \hyperref[def:absolute-error-approximation]{Absolute Error of an Approximation}
  \item \hyperref[thm:absolute-value-product]{Absolute Value of a Product}
  \item \hyperref[thm:absolute-value-sum-bound]{Generalised Triangle Inequality}
\end{itemize}
\end{dependencies}

\begin{theorembox}{Theorem (Absolute Error of a Quotient)}
\begin{theorem}[Absolute Error of a Quotient]
\label{thm:absolute-error-quotient}
Let $x,y,\widetilde{x},\widetilde{y}\in\mathbb{R}$ with
$y\neq 0$, $\widetilde{y}\neq 0$, and
$\delta(\widetilde{y})<1$. Then
\[
  \Delta\!\left(\frac{\widetilde{x}}{\widetilde{y}}\right)
  \leq
  \frac{
    |\widetilde{x}|\Delta(\widetilde{y})
    +
    |\widetilde{y}|\Delta(\widetilde{x})
  }{\widetilde{y}^{2}}
  \cdot
  \frac{1}{1-\delta(\widetilde{y})}.
\]

\smallskip\noindent
\hyperref[prf:absolute-error-quotient]{\textit{Go to proof.}}
\end{theorem}
\end{theorembox}

\begin{remark*}[Standard quantified statement]
\[
\forall x,y,\widetilde{x},\widetilde{y}\in\mathbb{R},\quad
\bigl(y\ne0\land \widetilde{y}\ne0\land \delta(\widetilde{y})<1\bigr)
\Rightarrow
\Delta\!\left(\frac{\widetilde{x}}{\widetilde{y}}\right)
\leq
\frac{|\widetilde{x}|\Delta(\widetilde{y})+|\widetilde{y}|\Delta(\widetilde{x})}{\widetilde{y}^{2}}
\cdot\frac{1}{1-\delta(\widetilde{y})}.
\]
\end{remark*}

\begin{remark*}[Predicate reading]
\[
\delta(\widetilde{y})<1\Rightarrow \Delta(\widetilde{x}/\widetilde{y})\leq \frac{|\widetilde{x}|\Delta(\widetilde{y})+|\widetilde{y}|\Delta(\widetilde{x})}{\widetilde{y}^{2}(1-\delta(\widetilde{y}))}.
\]
\end{remark*}

\begin{remark*}[Interpretation]
The factor $(1-\delta(\widetilde{y}))^{-1}$ records a real danger: division
is unstable when the denominator is poorly known relative to its own size.
This is why quotient estimates always require a hypothesis keeping the
denominator away from zero.
\end{remark*}

\begin{dependencies}
\begin{itemize}
  \item \hyperref[def:absolute-error-approximation]{Absolute Error of an Approximation}
  \item \hyperref[def:relative-error-approximation]{Relative Error of an Approximation}
  \item \hyperref[thm:absolute-value-quotient]{Absolute Value of a Quotient}
  \item \hyperref[thm:absolute-value-sum-bound]{Generalised Triangle Inequality}
\end{itemize}
\end{dependencies}

\begin{propositionbox}{Proposition (Relative Error of a Product)}
\begin{proposition}[Relative Error of a Product]
\label{prop:relative-error-product}
Let $\widetilde{x},\widetilde{y}\neq 0$. Then
\[
  \delta(\widetilde{x}\widetilde{y})
  \leq
  \delta(\widetilde{x})
  +
  \delta(\widetilde{y})
  +
  \delta(\widetilde{x})\delta(\widetilde{y}).
\]

\smallskip\noindent
\hyperref[prf:relative-error-product]{\textit{Go to proof.}}
\end{proposition}
\end{propositionbox}

\begin{remark*}[Standard quantified statement]
\[
\forall \widetilde{x},\widetilde{y}\in\mathbb{R}\setminus\{0\},\quad
\delta(\widetilde{x}\widetilde{y})
\leq \delta(\widetilde{x})+\delta(\widetilde{y})
+\delta(\widetilde{x})\delta(\widetilde{y}).
\]
\end{remark*}

\begin{remark*}[Predicate reading]
\[
\delta(\widetilde{x}\widetilde{y})\leq \delta(\widetilde{x})+\delta(\widetilde{y})+\delta(\widetilde{x})\delta(\widetilde{y}).
\]
\end{remark*}

\begin{remark*}[Interpretation]
Relative product error is approximately additive, with a second-order correction measuring simultaneous error in both factors.
\end{remark*}

\begin{dependencies}
\begin{itemize}
  \item \hyperref[def:relative-error-approximation]{Relative Error of an Approximation}
  \item \hyperref[thm:absolute-error-product]{Absolute Error of a Product}
  \item \hyperref[thm:absolute-value-product]{Absolute Value of a Product}
\end{itemize}
\end{dependencies}

\begin{propositionbox}{Proposition (Relative Error of a Quotient)}
\begin{proposition}[Relative Error of a Quotient]
\label{prop:relative-error-quotient}
Let $\widetilde{x},\widetilde{y}\neq 0$ and suppose
$\delta(\widetilde{y})<1$. Then
\[
  \delta\!\left(\frac{\widetilde{x}}{\widetilde{y}}\right)
  \leq
  \frac{\delta(\widetilde{x})+\delta(\widetilde{y})}
       {1-\delta(\widetilde{y})}.
\]

\smallskip\noindent
\hyperref[prf:relative-error-quotient]{\textit{Go to proof.}}
\end{proposition}
\end{propositionbox}

\begin{remark*}[Standard quantified statement]
\[
\forall \widetilde{x},\widetilde{y}\in\mathbb{R}\setminus\{0\},\quad
\delta(\widetilde{y})<1
\Rightarrow
\delta\!\left(\frac{\widetilde{x}}{\widetilde{y}}\right)
\leq
\frac{\delta(\widetilde{x})+\delta(\widetilde{y})}{1-\delta(\widetilde{y})}.
\]
\end{remark*}

\begin{remark*}[Predicate reading]
\[
\delta(\widetilde{y})<1\Rightarrow \delta(\widetilde{x}/\widetilde{y})\leq \frac{\delta(\widetilde{x})+\delta(\widetilde{y})}{1-\delta(\widetilde{y})}.
\]
\end{remark*}

\begin{remark*}[Interpretation]
Relative quotient error is product-like error amplified by uncertainty in the denominator.
\end{remark*}

\begin{dependencies}
\begin{itemize}
  \item \hyperref[def:relative-error-approximation]{Relative Error of an Approximation}
  \item \hyperref[thm:absolute-error-quotient]{Absolute Error of a Quotient}
  \item \hyperref[thm:absolute-value-quotient]{Absolute Value of a Quotient}
\end{itemize}
\end{dependencies}

\subsection*{Linearized Error Rules}

\begin{remark*}[Linearized error rules]
When the errors are small, second-order products such as
$\delta(\widetilde{x})\delta(\widetilde{y})$ are much smaller than the
first-order terms. One then informally reads
\[
  \delta(\widetilde{x}\widetilde{y})
  \approx
  \delta(\widetilde{x})+\delta(\widetilde{y}),
  \qquad
  \delta\!\left(\frac{\widetilde{x}}{\widetilde{y}}\right)
  \approx
  \delta(\widetilde{x})+\delta(\widetilde{y}).
\]
These are not replacement theorems for the exact estimates. They are
first-order approximations: useful for intuition, but justified only when
the discarded terms have been shown to be negligible at the relevant scale.
\end{remark*}

\begin{example*}[Catastrophic Cancellation]
Suppose two measured lengths are
\[
  H_1=200\pm 0.5\quad\text{and}\quad H_2=199.8\pm 0.5
\]
in centimeters. Their difference is approximately $0.2$, but the absolute
error bound is $0.5+0.5=1.0$. The relative error of the difference is large
because the denominator has collapsed. Subtraction of nearly equal quantities
can destroy relative precision even when each original measurement is quite
reasonable.
\end{example*}

\begin{example*}[From Product Error to Product Limits]
Let
\[
  x_n=2+\frac{1}{n},\qquad
  y_n=3-\frac{1}{n^2}.
\]
Write $x_n=2+\alpha_n$ and $y_n=3+\beta_n$, where
$\alpha_n=1/n$ and $\beta_n=-1/n^2$. Then
\[
  x_ny_n-6
  =
  2\beta_n+3\alpha_n+\alpha_n\beta_n.
\]
Each term on the right tends to zero. The product limit is therefore not a
new mystery; it is the finite product-error formula with the errors promoted
to null sequences.
\end{example*}
